We present an empirical study of how progressively adding task-specific guardrails to prompts impacts large language models’ ability to generate code optimization suggestions for HPC workloads. Across three prompt designs, we observe that tighter guardrails systematically shift model behavior from aggressive but low-utility transformations toward conservative, high-utility optimizations that compile, preserve numerical correctness, and yield measurable performance gains. The first prompt, which loosely describes the optimization task, produces many syntactically plausible but non-compiling or numerically incorrect suggestions and often requires up to ten optimization–feedback iterations to reach a viable configuration. The second prompt introduces explicit constraints on correctness, safety, and resource usage, reducing over-aggressive transformations and increasing the fraction of optimizations that lead to positive speedups without degrading numerical accuracy. Finally, our most constrained “agentic” prompt instantiates a closed-loop optimization routine in which the LLM autonomously proposes an optimization, checks compilation and numerical correctness against a reference, and revises its plan before returning a single final suggestion, achieving one-shot successful optimization on our benchmarks. We evaluate these prompts on diverse HPC kernels and systems, and we quantify their effects along three axes: performance gain, correctness, and optimization aggressiveness (e.g., code churn, magnitude of configuration change). Our results show that guardrail design is not merely a safety or policy concern, but a central lever for turning general-purpose code LLMs into reliable performance co-designers for scientific computing.